\subsection{Exponential Models Using Percents} 
Often, we describe multiplication in terms of \emph{percents} or \emph{percent change}.
\begin{Example}
Each year, approximately 65\% of all aluminum used is recycled. If 50\;kg (50,000 grams) of aluminum enter circulation this year, write an exponential function describing the amount $A$ of this aluminum (in grams) that remains in circulation after $t$ years.
\begin{itemize}
\item To find 65\% of some amount, we multiply by \makeblanksmall, so $m=\makeblanksmall$
\item We start with \makeblank[1in] grams of aluminum, so $A_0=\makeblank[1in]$
\end{itemize}
\[
A(t)=\makeblank
\]
\end{Example}
\begin{Example}
The population of a small town decreases by 4\% each year. This year, the population of the town is 13,850. Write an exponential function describing the population $p$ of the town $t$ years from now.
\begin{itemize}
\item If the population decreases by 4\%, then \makeblank remains. So each year we multiply by \makeblank. So $m=\makeblanksmall$.
\item The starting population is \makeblank[1in], so $p_0=\makeblank[1in]$.
\end{itemize}
\[
p(t)=\makeblank
\]
\end{Example}
\begin{Example}
The value of a certain stock increases by 12\% each year. This year, the stock is worth \$350 per share. Write an exponential function describing the value $v$ of the stock $t$ years from now.
\begin{itemize}
\item If the value increases by 12\%, it ends up with \makeblank of its original value. So each year we multiply by \makeblank. So $m=\makeblanksmall$.
\item The starting value is \makeblank[1in], so $v_0=\makeblanksmall$
\end{itemize}
\[
v(t)=\makeblank
\]
\end{Example}
\noindent When we have a percent change:
\begin{itemize}
\item If we have a percent \emph{increase} we \makeblank.
\item If we have a percent \emph{decrease} we \makeblank.
\end{itemize}
